%! AP Statistics — Unit 1, Lesson 1.6 — Homework ANSWER KEY
\documentclass[10pt]{article}
\usepackage{apstats-article}
\usepackage{apstats-key}

%=============================================================================
\begin{document}

\pageheader{Unit 1, Lesson 1.6}{Homework --- Answer Key}

\vspace{0.1in}

\begin{tcolorbox}[colback=sky, colframe=navy, boxrule=0.8pt, arc=2mm,
    left=3mm, right=3mm, top=1.5mm, bottom=1.5mm]
\small For all written descriptions: use the SOCS framework; include context
(variable name and population) in every sentence; use complete sentences.
AP free-response graders look for all four components and context throughout.
\end{tcolorbox}

\vspace{0.14in}

% ════════════════════════════════════════════════════════════════════════════
% PART A
% ════════════════════════════════════════════════════════════════════════════
\boxguard
\begin{blurbbox}[Part A \quad Identify the Shape]{navylight}
\small

For each distribution described below, identify the most likely shape from
this list: \textit{symmetric, skewed right, skewed left, bimodal, uniform}.
Give one piece of evidence from the description that supports your answer.

\vspace{10pt}

\begin{enumerate}[leftmargin=1.5em, itemsep=14pt]

  \item The prices of homes sold in a county last year. Most homes sold for
        \$180{,}000--\$350{,}000, but a small number of luxury estates sold
        for over \$2 million.\\[4pt]
        Shape: \ans{Skewed right} \quad Evidence: \ans{long right tail}\\[1pt]
        \ansline{Most homes sell in the \$180k--\$350k bulk; estates above \$2M pull the tail right.}

  \item Scores on a very easy quiz (out of 10) for a class of 30 students.
        The distribution shows most students scored 9 or 10, with only a few
        scoring below 7.\\[4pt]
        Shape: \ans{Skewed left} \quad Evidence: \ans{tail toward low scores}\\[1pt]
        \ansline{Most students scored 9--10 (right bulk); a few below 7 form a left tail.}

  \item The number of hours of sleep per night for 60 adults, with a peak
        near 7 hours and roughly equal drop-off on both sides.\\[4pt]
        Shape: \ans{Symmetric, unimodal} \quad Evidence: \ans{equal drop-off}\\[1pt]
        \ansline{The peak near 7 hours falls off about equally on both sides --- a bell shape.}

  \item Survey responses on a 5-point scale where all responses (1, 2, 3, 4,
        and 5) were chosen by nearly equal numbers of respondents.\\[4pt]
        Shape: \ans{Uniform} \quad Evidence: \ans{bars all equal}\\[1pt]
        \ansline{All five options drew nearly equal counts, so every bar is the same height.}

\end{enumerate}

\end{blurbbox}

\vspace{0.14in}

% ════════════════════════════════════════════════════════════════════════════
% PART B
% ════════════════════════════════════════════════════════════════════════════
\boxguard
\begin{blurbbox}[Part B \quad Write a Complete SOCS Description]{greenacc}
\small

A researcher records the \textbf{number of minutes per day} that 45 middle
school students spend on physical activity. Key features from the histogram:

\begin{itemize}[itemsep=3pt]
  \item Shape: skewed right; most students are active 20--60 minutes per day.
  \item A long tail extends to the right: a few students report 90--120 minutes.
  \item One student reported 150 minutes --- far above the rest of the group.
  \item The center of the distribution is approximately 38 minutes.
  \item Values range from 5 to 150 minutes.
\end{itemize}

\vspace{10pt}

\begin{enumerate}[leftmargin=1.5em, itemsep=10pt]

  \item Write a complete SOCS description of the distribution of daily physical
        activity time for these 45 middle school students. Use complete sentences
        with context in every sentence. Label each sentence S, O, C, or S in
        the margin.\\[6pt]
        \ansline{[S] The distribution of daily physical activity time for these 45 middle}\\[1pt]
        \ansline{school students is skewed right, with most active 20--60 minutes a day.}\\[1pt]
        \ansline{[O] One student reported 150 minutes, a potential outlier far above the rest.}\\[1pt]
        \ansline{[C] The center is about 38 minutes per day --- a typical activity level for}\\[1pt]
        \ansline{a middle school student in this sample.}\\[1pt]
        \ansline{[S] Times range from 5 to 150 minutes, a spread of 145 minutes.}

  \item If the one student reporting 150 minutes is removed, predict how (if
        at all) the shape, center, and spread of the distribution would change.
        Explain your reasoning.\\[6pt]
        \ansline{\textit{Shape:} still skewed right but less so --- the long tail shortens.}\\[1pt]
        \ansline{\textit{Center:} drops slightly; the mean falls more than the median.}\\[1pt]
        \ansline{\textit{Spread:} drops noticeably --- the max falls from 150 to roughly 90--120.}

  \item Based on the shape of this distribution, would you expect the mean
        daily activity time to be greater than, equal to, or less than the
        median? Explain why.\\[6pt]
        \ansline{\textbf{Greater than the median.} In a right-skewed distribution the few}\\[1pt]
        \ansline{large values pull the mean upward more than they move the median, so the}\\[1pt]
        \ansline{mean exceeds the approximate median of 38 minutes.}

\end{enumerate}

\end{blurbbox}

\vspace{0.14in}

% ════════════════════════════════════════════════════════════════════════════
% PART C
% ════════════════════════════════════════════════════════════════════════════
\boxguard
\begin{blurbbox}[Part C \quad AP-Style Free Response]{redacc}
\small

The following summary describes the distribution of \textbf{fuel efficiency
(miles per gallon)} for a random sample of 60 cars tested by a consumer group.

\vspace{8pt}
\renewcommand{\arraystretch}{1.5}
\begin{tabularx}{0.72\linewidth}{@{} >{\bfseries}p{0.30\linewidth}
                                      X @{}}
\rowcolor{sky}
\textbf{Feature} & \textbf{Value / Description} \\
\midrule
Shape & Roughly symmetric, unimodal \\
Minimum & 18 mpg \\
Maximum & 47 mpg \\
Approximate center & 32 mpg \\
Outliers & One car at 47 mpg; no low-end outliers \\
\end{tabularx}

\vspace{10pt}

\begin{enumerate}[leftmargin=1.5em, itemsep=10pt]

  \item Describe the distribution of fuel efficiency for these 60 cars using
        the SOCS framework. Write in complete sentences and include context in
        every sentence.\\[6pt]
        \ansline{[S] Fuel efficiency for these 60 cars is roughly symmetric and unimodal.}\\[1pt]
        \ansline{[O] One car at 47 mpg is a potential high outlier; none at the low end.}\\[1pt]
        \ansline{[C] The center is about 32 mpg --- typical efficiency for a car in this sample.}\\[1pt]
        \ansline{[S] Values run from 18 to 47 mpg, a spread of 29 mpg.}

  \item The consumer group adds 5 more cars to the sample; all 5 have fuel
        efficiency of 14 mpg. Describe how each SOCS component would change.\\[6pt]
        \ansline{\textit{Shape:} shifts toward skewed left --- the new cars build a low tail.}\\[1pt]
        \ansline{\textit{Outliers:} the five 14-mpg cars become low outliers, below the old min.}\\[1pt]
        \ansline{\textit{Center:} decreases; the mean falls more than the median.}\\[1pt]
        \ansline{\textit{Spread:} increases --- the minimum drops from 18 to 14 mpg.}

  \item A classmate says that for a symmetric distribution the mean and median
        should be about the same. Is the classmate correct? Explain.\\[6pt]
        \ansline{Yes, for a roughly symmetric distribution. Values on either side of}\\[1pt]
        \ansline{center offset one another, so the mean is not pulled either way. The}\\[1pt]
        \ansline{47-mpg outlier lifts the mean slightly, but the gap stays small.}

\end{enumerate}

\end{blurbbox}

\vspace{0.12in}

\boxguard
\begin{extensionbox}
\small
\textbf{Optional Extension --- Real Data SOCS Description}\\[4pt]
Find a real histogram, dotplot, or stemplot in a news article, textbook, or
data website (e.g., census.gov, gapminder.org). Write a complete SOCS
description of the distribution. Then explain: what does the shape of this
distribution suggest about the real-world process that generated it?

\vspace{8pt}
\writeline\\[1pt]\writeline\\[1pt]\writeline\\[1pt]\writeline\\[1pt]\writeline

\vspace{4pt}
\textcolor{slate}{\scriptsize Source / URL:\;\blank{11cm}}
\end{extensionbox}

\end{document}
